\section{Significance of the Study}

Graph-based representations allow for a novel analysis of music pieces as it allows for the relationships between segments throughout the piece to be captured. The structure of music pieces may go beyond the individual music events as a holistic, nonlinear structure is formed among the events \cite{HernandezOlivanBeltran}. 

Reducing the runtime of this model will allow for increased speed of analysis of music pieces as it will require less resources. This will allow for more comparisons, allowing for improved analysis of pieces as there can be a larger number of comparisons that can be made and increasing the value of each individual comparison \cite{ClarkCook}. Furthermore, this reduces the need of the model on modules in creating checkpoints, such as pickle, thus reducing the potential risk of data loss and wasted computation time. A faster implementation of the model also increases its viability for industry-level applications. 

In addition to the improved similarity comparisons through the graph-representations, DTW is a popular algorithm for many purposes. Exploring a faster implementation of this algorithm and its limitations can improve applications of DTW in many other models, not limited to computational musicology. Furthermore, it allows for further research on DTW and its applications.